% from ega4-21b-fr.tex
\oldpage[IV-4]{276}
Montrons d'abord que l'homomorphisme $\varinjlim\operatorname{Div}(U_\lambda)\to\mathfrak{Z}^1(X)$ est \emph{injectif}.
Posons $T=\bigcap_\lambda U_\lambda$, et notons que les $U_\lambda$ forment un \emph{système fondamental de voisinages} de $T$.
En effet, pour tout ouvert $V\supset T$, $X-V$ est un sous-espace fermé de $X$, donc un espace noethérien dont toute partie fermée irréductible admet un point générique;
comme les ensembles $(X-V)\cap(X-U_\lambda)$ sont fermés dans $X-V$, forment une famille filtrante croissante et ont pour réunion $X-V$, notre assertion résulte de \sref[0]{0.9.2.4}.
Cela étant, il s'agit de voir que si $D\in\operatorname{Div}(U_\lambda)$ est tel que $\operatorname{cyc}(D)=0$ dans $\mathfrak{Z}^1(U_\lambda)$, alors il existe $\mu\geq\lambda$ tel que $D|U_\mu=0$.
En vertu de ce qui précède, il suffira de voir que pour tout $x\in T$, on a $D_x=0$;
en effet, il y aura alors un voisinage $W(x)$ de $x$ dans $X$ tel que $D|W(x)=0$, et la réunion des $W(x)$ pour $x\in T$ contient un $U_\mu$.
Compte tenu de \sref{IV.21.4.6}, on est donc ramené au cas où $X=\Spec(\sh{O}_{X,x})$ avec $x\in T$;
mais par hypothèse $\sh{O}_{X,x}$ est factoriel, donc intègre et intégralement clos, et $\Spec(\sh{O}_{X,x})$ est alors normal;
donc la conclusion résulte de \sref{IV.21.6.9}, (i).

